\section{Architecture and design}
\label{sec:design}

\subsection{Overall architecture}

NEO-Watch is made of three parts, shown in Figure~\ref{fig:architecture}:
\begin{itemize}
  \item a \textbf{REST web service} written with FastAPI (the backend), which holds all the
    business logic and is the only part that reads and writes the database;
  \item a \textbf{user interface} written with Streamlit (the frontend), which only calls the
    web service over HTTP and never accesses the database directly;
  \item a \textbf{PostgreSQL database}, filled from the NASA NeoWs API.
\end{itemize}

\begin{figure}[H]
  \centering
  \begin{tikzpicture}[
      box/.style={draw, rounded corners, align=center, minimum width=3.4cm, minimum height=1cm, font=\small},
      layer/.style={draw, align=center, minimum width=5.2cm, minimum height=0.75cm, font=\small},
      arrow/.style={-{Stealth}, thick},
      node distance=0.15cm]
    \node[box] (front) {Streamlit frontend\\ \footnotesize pages, API client};
    \node[layer, right=1.8cm of front, yshift=1.5cm] (ctrl) {Controllers (FastAPI routers)};
    \node[layer, below=of ctrl] (svc) {Services (business logic)};
    \node[layer, below=of svc] (dao) {DAO (SQL queries)};
    \node[draw, dashed, fit=(ctrl)(svc)(dao), inner sep=0.25cm, label={above:\small FastAPI backend}] (back) {};
    \node[box, below=1.1cm of dao] (db) {PostgreSQL};
    \node[box, right=1.2cm of svc] (nasa) {NASA APIs\\ \footnotesize \texttt{/feed}, \texttt{/apod}};
    \draw[arrow] (front.east) -- node[above, font=\footnotesize, align=center]{HTTP / JSON\\ JWT token} (back.west);
    \draw[arrow] (dao) -- (db);
    \draw[arrow] (svc) -- (nasa);
  \end{tikzpicture}
  \caption{Overall architecture of NEO-Watch}
  \label{fig:architecture}
\end{figure}

This separation has three advantages. The database is protected: every access goes through the
web service, which checks the authentication and the business rules. The interface can be
replaced without touching the backend. Finally, the web service can be explored and tested on
its own, through the documentation generated by FastAPI (Swagger, at \texttt{/docs}).

\subsection{A layered backend}

The backend follows the layered architecture of the course: each layer only calls the layer
below it. Table~\ref{tab:layers} gives the role of each package.

\begin{table}[H]
  \centering
  \small
  \begin{tabularx}{\textwidth}{@{}l l>{\raggedright\arraybackslash}X@{}}
    \toprule
    \textbf{Layer} & \textbf{Package} & \textbf{Role} \\
    \midrule
    Controller & \texttt{controller} & HTTP endpoints: reads the request, checks the token and the
      rights, turns errors into HTTP status codes. No SQL, no business rule. \\
    Schema & \texttt{schema} & Pydantic models of the requests and responses: types and bounds
      are checked automatically (error 422 otherwise). \\
    Service & \texttt{service} & Business rules (for instance: a NEO needs a name and positive
      diameters), NASA calls, import scheduling. \\
    DAO & \texttt{dao} & One class per table or feature; the only place where SQL is written.
      Rows are turned into business objects. \\
    Business object & \texttt{business\_object} & Domain classes (\texttt{Neo}, \texttt{Approach},
      \texttt{User}, \texttt{Alert}\ldots) with private attributes and properties. \\
    Utils & \texttt{utils} & Security (passwords, JWT), logs, singleton, database reset. \\
    \bottomrule
  \end{tabularx}
  \caption{Packages of the backend}
  \label{tab:layers}
\end{table}

Thanks to this organisation, each layer can be tested alone: the service tests replace the DAOs
with mocks, and the DAO tests run on a dedicated database schema (Section~5).

\subsection{Use cases}

Figure~\ref{fig:usecase} shows the use cases of the two actors. An administrator is a user with
more rights: they can do everything a user does, and also trigger the NASA import, manage the
accounts and reset the database.

\begin{figure}[H]
  \centering
  \includegraphics[width=0.62\textwidth]{figures/use_case_diagram.png}
  \caption{Use case diagram}
  \label{fig:usecase}
\end{figure}

\subsection{Business objects}

Figure~\ref{fig:classes} shows the business classes. A \texttt{Neo} has several
\texttt{Approach} objects (one per day when it comes close to a celestial body). A \texttt{User}
has favourites, alerts, connections and imports. Every attribute is private and exposed through
properties, so that a value can be checked or converted when it is changed. Each class has a
\texttt{to\_dict} method, used to build the JSON answers of the web service.

\begin{figure}[H]
  \centering
  \includegraphics[width=0.95\textwidth]{figures/class_diagram.png}
  \caption{Class diagram of the business objects}
  \label{fig:classes}
\end{figure}

Each service uses one or several DAOs: most services use the DAO of their feature;
\texttt{ImportService} writes NEOs, approaches and imports, and \texttt{FavouriteService} reads
the favourites and the approaches.

% TODO: the service/DAO diagram lists methods that do not exist in the code
% (create_user_neo, update_user_neo, delete_user_neo...). Put it back once it is fixed,
% and add "Figure~\ref{fig:servicedao} shows..." to the paragraph above.
% \begin{figure}[H]
%   \centering
%   \includegraphics[width=0.95\textwidth]{figures/service_dao_diagram.png}
%   \caption{Services and the DAOs they use}
%   \label{fig:servicedao}
% \end{figure}

\subsection{Data model}

Figure~\ref{fig:pdm} shows the physical data model. The main design choices are the following.
\begin{itemize}
  \item \textbf{\texttt{neo} and \texttt{approach} are two tables.} An object passes several
    times; its physical characteristics are stored once, its approaches as many times as needed.
    The constraint \texttt{UNIQUE (neo\_id, approach\_day)} allows at most one approach per object
    and per day.
  \item \textbf{Objects created by users are protected.} The \texttt{origin} column is
    \texttt{NASA} or \texttt{USER}, and \texttt{nasa\_id} is empty for user objects. The import
    updates an existing NASA object when its \texttt{nasa\_id} is already known
    (\texttt{INSERT \ldots ON CONFLICT (nasa\_id) DO UPDATE}); a user object can never match.
  \item \textbf{Imports are idempotent.} Since NEOs and approaches are inserted or updated
    (\emph{upsert}), importing the same period twice creates no duplicate.
  \item \textbf{Deletions cascade.} Deleting an account deletes its favourites, alerts and
    connections (\texttt{ON DELETE CASCADE}); the NEOs declared by the user are kept, as decided
    in the analysis report.
  \item \textbf{The \texttt{import} table is a log.} Each import records its period, its trigger
    (first import, update or daily refresh), its status (\texttt{running}, \texttt{success},
    \texttt{partial}, \texttt{failed}) and its counters. The next import uses it to know where to
    resume.
  \item \textbf{Indexes} are created on the foreign keys and on the approach day, which is the
    most frequent filter (dashboard, search by date, alerts).
\end{itemize}
The \texttt{search} and \texttt{proposal} tables come from the analysis report and are not used
by the final application (Table~\ref{tab:changes}).

\begin{figure}[H]
  \centering
  \includegraphics[width=\textwidth]{figures/physical_data_model.png}
  \caption{Physical data model}
  \label{fig:pdm}
\end{figure}

\subsection{Web service}

Table~\ref{tab:endpoints} lists the endpoints of the web service. All of them, except logging in
and creating an account, require an access token in the \texttt{Authorization} header. The
endpoints marked \emph{admin} are refused to other users (error 403), and a user can only read or
change their own favourites and alerts.

{\small
\begin{longtable}{@{}l l p{5.2cm}@{}}
  \toprule
  \textbf{Method} & \textbf{Path} & \textbf{Description} \\
  \midrule
  \endfirsthead
  \toprule
  \textbf{Method} & \textbf{Path} & \textbf{Description} \\
  \midrule
  \endhead
  \bottomrule
  \caption{Endpoints of the web service}
  \label{tab:endpoints}
  \endlastfoot
  POST & \texttt{/login/} & Log in; returns an access token and a refresh token \\
  POST & \texttt{/login/token} & Same, with a form (used by the Swagger \emph{Authorize} button) \\
  POST & \texttt{/login/register} & Create an account \\
  POST & \texttt{/login/refresh} & New tokens from a valid refresh token \\
  GET & \texttt{/login/me} & Profile of the logged-in user \\
  GET & \texttt{/users/} & List the accounts (admin) \\
  PUT & \texttt{/users/\{user\_id\}/admin} & Give the administrator role (admin) \\
  DELETE & \texttt{/users/\{user\_id\}} & Delete an account (admin) \\
  GET & \texttt{/users/\{user\_id\}/connections} & Connection history of a user (admin) \\
  GET & \texttt{/dashboard/} & Dashboard of a period (default: next 7 days) \\
  GET & \texttt{/neo} & List the NEOs, with pagination \\
  POST & \texttt{/neo/search} & Multi-criteria search, sorted and paginated \\
  GET & \texttt{/neo/\{neo\_id\}} & Details of a NEO and its approaches \\
  POST & \texttt{/neo} & Declare a NEO (origin \texttt{USER}) \\
  POST & \texttt{/favourites/} & Add a favourite \\
  DELETE & \texttt{/favourites/\{user\_id\}/\{neo\_id\}} & Remove a favourite \\
  GET & \texttt{/favourites/user/\{user\_id\}} & Favourites of a user \\
  GET & \texttt{/favourites/user/\{user\_id\}/approaches} & Approaches of the favourites \\
  GET & \texttt{/favourites/user/\{user\_id\}/neo/\{neo\_id\}/graph} & Distance plot of a favourite (PNG) \\
  POST & \texttt{/alerts/} & Create an alert \\
  GET & \texttt{/alerts/user/\{user\_id\}} & Alerts of a user \\
  GET & \texttt{/alerts/user/\{user\_id\}/notifications} & New notifications of a user \\
  PUT & \texttt{/alerts/\{alert\_id\}/toggle} & Switch an alert on or off \\
  DELETE & \texttt{/alerts/\{alert\_id\}} & Delete an alert \\
  GET & \texttt{/statistics/} & Statistics on the NEOs \\
  GET & \texttt{/apod/} & Astronomy Picture of the Day (optional date) \\
  POST & \texttt{/import/update} & Start a NASA update in the background (admin) \\
  GET & \texttt{/import/status} & Is an import running? \\
  GET & \texttt{/import/last} & End date of the last import \\
  POST & \texttt{/reset\_database} & Re-create the tables and the seed data (admin, refused during an import) \\
\end{longtable}
}

Figure~\ref{fig:sequence} follows a search request through the layers.

\begin{figure}[H]
  \centering
  \small
  \begin{tabularx}{\textwidth}{@{}r>{\raggedright\arraybackslash}X@{}}
    1. & The Streamlit page sends \texttt{POST /neo/search} with the criteria in JSON and the token. \\
    2. & FastAPI checks the token (\texttt{get\_current\_user}) and validates the criteria (Pydantic). \\
    3. & \texttt{neo\_controller} calls \texttt{NeoService.search}, which calls \texttt{NeoDao.find}. \\
    4. & \texttt{NeoDao} builds a parameterised SQL query, joins \texttt{approach} only for a date criterion. \\
    5. & PostgreSQL returns the rows; the DAO turns them into \texttt{Neo} objects. \\
    6. & The controller returns them as JSON; the page displays a table and offers a CSV export. \\
  \end{tabularx}
  \caption{Path of a search request through the layers}
  \label{fig:sequence}
\end{figure}

\subsection{User interface}

The interface is a multi-page Streamlit application (Table~\ref{tab:pages}). After logging in,
the user reaches a menu that gives access to every feature; administrators see an extra section
with the import and account tools. The new alert notifications are shown on this menu as soon as
the user logs in.

\begin{table}[H]
  \centering
  \small
  \begin{tabularx}{\textwidth}{@{}l>{\raggedright\arraybackslash}X@{}}
    \toprule
    \textbf{Page} & \textbf{Content} \\
    \midrule
    \texttt{home.py} & Log in and create an account \\
    \texttt{neo\_menu.py} & Menu, new notifications, administrator tools (manual update, reset) \\
    \texttt{dashboard.py} & Dashboard of the chosen period (F2) \\
    \texttt{search\_neo.py} & Search, details of a NEO, CSV export, add to favourites (F3) \\
    \texttt{create\_neo.py} & Declare a NEO (F4) \\
    \texttt{favourites.py} & Favourites, approaches, distance plot, CSV and PNG export (F5) \\
    \texttt{alert.py} & Alerts and notifications (F6) \\
    \texttt{statistics.py} & Statistics and charts (FO3) \\
    \texttt{apod.py} & Astronomy Picture of the Day (FO8) \\
    \texttt{admin.py} & Accounts and connection history, for administrators (F1) \\
    \bottomrule
  \end{tabularx}
  \caption{Pages of the Streamlit interface}
  \label{tab:pages}
\end{table}

Two utilities are shared by all the pages. \texttt{api\_client.py} adds the token to every
request; when the access token has expired (it lasts 30 minutes), it asks for a new one with the
refresh token and sends the request again, so that the user is not logged out in the middle of
their work. \texttt{auth\_guard.py} sends a user who is not logged in back to the login page.
